\chapter{The free-form spike: a shape basis and a gradient}
\label{ch:spline}

Every plug computed so far in this document inherited its shape. The
spike was a \emph{streamline} of the expansion fan at the cowl lip:
you start at a radius and integrate the flow direction downstream, and
the contour is whatever comes out. Nothing was ever chosen. The only
two knobs were the angle at which the flow arrives
(Chapter~\ref{ch:compare}) and the station at which the spike is cut
(Chapter~\ref{ch:plugcycle}).

This chapter gives the spike a \emph{shape}: a family of contours with
free parameters, a derivative of thrust with respect to each of them,
and an optimizer that moves them. It is the last structural piece of
the program, and it is what finally makes the plug and the bell
comparable at equal freedom.

\section{Why the inherited contour is not the answer}
\label{sec:whynotstreamline}

The streamline construction is not arbitrary --- it solves a real
problem, exactly. Take the centred expansion fan that springs from the
cowl lip and let the spike follow a streamline through it: every
characteristic completes its turning, the flow leaves axially, and the
exit pressure lands exactly on ambient. That is the \emph{ideal},
full-length plug, and for that problem the streamline is optimal.

Then the spike is cut at a finite length. The contour was shaped to
deliver a downstream state that is now never reached, and nothing in
its construction knows the cut exists. The criterion that produced it
is no longer the criterion being maximized.

What the fixed-length problem rewards is visible in the thrust
integral itself. The axial force the spike collects is
\begin{equation}
  J_{\text{wall}} \;=\; \int (p - p_a)\, 2\pi y \,(-\mathrm{d}y),
  \label{eq:wallpush}
\end{equation}
and the weight is $2\pi y$: a given overpressure is worth far more
where the spike is still \emph{fat} than near its tip, simply because
there is more area there for it to act on. Under a length constraint
the right move is therefore to bias the expansion --- hold the
pressure up early, at large radius, and accept a worse state at the
end that is about to be discarded anyway. The streamline cannot do
this. It distributes its turning as though the whole contour mattered
equally, because it is not choosing: it is being traced.

This is Rao's argument for the bell (Chapter~\ref{ch:referee}),
transposed. His problem \emph{is} the length-constrained one, and the
multiplier $\lambda_3$ that Section~\ref{sec:o33} certified against
the automatic adjoint is precisely the shadow price of that length
constraint. The plug carries the same variational structure. What was
missing was not the theory but the machinery to act on it.

\section{The basis is borrowed, deliberately}

The wall is a \textbf{clamped-left natural cubic spline}: a curve
assembled from cubic pieces joined so that value, slope and curvature
are continuous across every join, pinned at its left end to a
prescribed height and slope. Its free parameters are the heights
$y_1, \dots, y_m$ at $m$ fixed stations, called \emph{knots}.

This is not a new basis written for the plug. It is the same routine
the bell optimizer uses, called directly. That choice is deliberate on
three counts. A bell-versus-plug comparison is only honest if both
sides have the \emph{same} shape freedom, which requires the same
basis. Any property already established for that basis --- in
particular that adding knots eventually represents any smooth contour
--- transfers without re-argument. And the adaptive-knot machinery
built for the bell, which decides \emph{where} knots should go rather
than accepting them evenly spaced, applies to this wall unchanged.
Section~\ref{sec:reprladder} shows that the plug needs it more than
the bell does.

\section{What is free, and what is not}

The design vector is
\begin{equation}
  W \;=\; [\,y_1,\; y_2,\; \dots,\; y_m\,],
\end{equation}
the spike radius at $m$ frozen stations spanning the contour, with the
last knot sitting exactly at the design length $L$.

Three quantities are \emph{not} free, and each for a reason.

\emph{The length} is fixed by placing the final knot at $L$. Without
it the optimization has a trivial answer: a longer spike always
collects more, so the optimizer would simply run away. Fixing the
length is what makes the problem well posed, and it is the same
constraint Rao imposed.

\emph{The mass flow} is fixed by construction. The starting radius of
the spike is found by a root search so that the flow crossing the
inlet station carries the reference mass flow. That station lies
\emph{upstream} of every knot, and the data on it are built from the
fan alone, so no shape parameter can reach it: the mass constraint and
the shape parameters are exactly independent. This is a structural
fact about the construction, not a numerical coincidence, and it means
the optimizer never has to trade thrust against mass.

\emph{The slope at the start} is fixed by physics. A wall is a
streamline, so where the incoming flow meets it the wall must run
along the flow. The clamped left slope is therefore the local flow
angle, not a design choice.

The inlet angle itself --- the inclination of the straight line by
which the engine delivers its flow --- remains the outer variable that
Chapter~\ref{ch:compare} swept, and it is not re-optimized here. The
two freedoms are kept apart so that a gain in thrust is attributable
to one of them rather than to their sum.

\section{The gradient: an adjoint through a frozen schedule}

The optimizer needs $\partial J/\partial y_k$ for every knot. Finite
differences would cost two marches per knot and would inherit the
step-size dilemma; the program has a better instrument.

The march makes \emph{decisions} as it goes --- which stations bracket
a wall point, how the cells are wired together. Those decisions are
discrete, and a derivative cannot be taken through them. The device,
used already by the bell optimizer, is to \textbf{record} them: one
march runs with concrete arithmetic and writes down every decision it
made, and every later evaluation \emph{replays} that record, treating
the decisions as frozen constants while the arithmetic stays
differentiable. The result is the exact derivative of the flow
\emph{at fixed topology}, obtained in one reverse pass regardless of
how many knots there are.

The record is refreshed whenever the optimizer accepts a step, so the
frozen topology is never allowed to drift far from the true one. The
agreement between this adjoint and central finite differences is
measured at every run and is reported in Section~\ref{sec:splineres}.

\section{The driver: trust-region SQP, in full}
\label{sec:splinedriver}

The optimizer deserves a complete description, both because every
number in this chapter and the next was produced by it and because
three of its defects --- each yielding plausible output while being
wrong --- are lessons in themselves.

\paragraph{What SQP is.} Sequential quadratic programming attacks a
nonlinear optimization by solving a sequence of tractable stand-ins:
at the current design $W$ it builds a \emph{quadratic model} of the
objective --- the measured gradient plus a curvature matrix
accumulated from successive gradient differences (BFGS: the model
learns curvature as it walks, no second derivatives are ever
computed) --- and, when constraints are present, a linear model of
those; the model problem is solved exactly, the true objective is
evaluated at the proposed step, and the cycle repeats. In the
\emph{trust-region} variant the model problem is solved only within
a radius $\Delta$ around $W$ --- the region where the model is
trusted --- and the ratio of \emph{actual} to \emph{predicted}
improvement governs $\Delta$: a step that delivers what the model
promised grows the radius, a step that does not shrinks it and is
rejected. This is the standard machinery for smooth nonconvex
problems, and the implementation underneath is scipy's
\code{trust-constr}.

\paragraph{Why the constrained problem arrives unconstrained.} The
three constraints of Section~15.3 --- mass, length, tangency ---
are imposed \emph{by construction} (the start radius by root
search, the last knot pinned at $L$, the clamped slope), so the SQP
never sees them: what remains is bound-free maximization over the
knot heights, and the sequential-quadratic machinery reduces to
trust-region quasi-Newton. The structure that makes the driver
nontrivial is not the constraints but the \emph{objective}: $J$ is
differentiable only through a frozen march record
(Section~15.4), and the record is itself position-dependent.

\paragraph{Decisions, records, and replays --- the vocabulary,
concretely.} Three words carry this section and deserve exact
meanings.

A \emph{decision} is any discrete choice the march makes while
building the field --- an integer or a branch, not a number that
varies smoothly with the design. The march makes thousands: which
segment of the previous column \emph{brackets} the characteristic
arriving at a new wall point (the wall-foot search returns an index,
and stepping columns further back changes that index by whole
units); how many rows a given column carries (the mesh gains a row
per column as the jet widens, and \emph{when} it gains one is a
threshold test); which cells are wired to which neighbours; which
side of every safety guard the arithmetic lands on. Each decision is
a step function of the knot heights: nudge the wall and the bracket
index, say, stays put --- until some threshold is crossed and it
jumps by one. A derivative cannot pass through a jump, so the march
as written is not differentiable, even though every formula inside
it is.

The \emph{record} (or \emph{schedule}) is the cure: one march runs
with concrete arithmetic and writes down every decision it made ---
the bracket indices, the row counts, the wiring, the branch
outcomes --- as plain integers in a ledger. Nothing about the flow
is stored; only the choices.

A \emph{replay} re-runs the march with every decision \emph{read
from the ledger instead of recomputed}. The comparisons and searches
vanish; what remains is a fixed pipeline of smooth
arithmetic --- the same cells, wired the same way, solving the same
equations --- through which the knot heights flow. That pipeline
\emph{is} differentiable, and one reverse pass through it yields
the exact gradient of $J$ \emph{at frozen topology}: the derivative
of the flow the march would compute \emph{if it kept making the
same choices}.

The price is locality, and it is the price the whole driver is
shaped around. Near the recorded base the frozen choices are the
right ones and the replayed $J$ agrees with the true $J$ to
round-off. Move the design far enough that the true march would
have decided differently --- a bracket shifts, a column gains its
row earlier --- and the replay, still faithfully executing the old
choices, computes a flow the true march would never produce: the
replayed objective is only \emph{piecewise} smooth, exact within
each decision cell and wrong across its boundary. The fold in
Figure~\ref{fig:s23driver}(b) is such a boundary seen edge-on, and
the finite-difference noise that once condemned a correct adjoint
(the ladder lesson of Section~15.6) is the same boundaries seen at
step-size scale.

\paragraph{The segmented walk.} Everything above dictates the
driver's shape: walk only on a record, distrust the walk the moment
it may have out-run the record's validity, and refresh the record
whenever a step is worth keeping. The driver therefore proceeds in
\emph{segments}:

\begin{enumerate}
\item \textbf{Record} at the current base $W$: one concrete march,
  writing down every discrete decision. The record must be
  \emph{Newton-certified} (every cell converged to round-off); an
  uncertified record rejects the base outright.
\item \textbf{Walk on the record}: one \code{trust-constr} call,
  limited to a few model iterations, maximizes the \emph{replayed}
  $J$ --- exact value and exact adjoint gradient at frozen topology
  --- inside the driver's own radius.
\item \textbf{Treat the endpoint as a trial, not an answer.} The
  fresh record at the endpoint must certify, and the replayed $J$
  there must beat the incumbent. Acceptance re-records and grows
  the radius ($\times 1.5$, capped); rejection reverts to the best
  certified design and halves the radius toward a floor.
\item \textbf{Stop} when a trial is worse at the radius floor, when
  the walk stops moving at the floor, or at the segment budget. The
  answer is the \emph{best certified design ever seen}, never the
  last endpoint.
\end{enumerate}

Figure~\ref{fig:s23driver}(a) shows a production walk in this
discipline --- fifteen accepted bases, three trials rejected for
thrust, two bases rejected by the certification gate, five
no-motion retries --- worth a thousand words of pseudocode.

\begin{figure}[htbp]
\centering
\includegraphics[width=\textwidth]{fig_s23_driver.pdf}
\caption{The driver at work, from the run of record. (a)~The
segmented trust-region SQP re-optimizing the spike at the doubled
resolution: the staircase is the accepted thrust (relative to the
start, in kN); red arrows mark trials the acceptance test rejected
(their thrust printed below, in kN --- the worst overshoot is
$-687$ kN, which the old driver would have adopted); purple crosses
mark bases rejected by the per-cell Newton certification gate;
orange triangles mark no-motion retries, where the inner solver
returned its start point and the outer driver halved the radius and
tried again on the same record. (b)~The mechanism behind those
retries, measured: thrust along the gradient direction is smooth and
matches the linear prediction (dashed) over five decades of step
size, then \emph{folds} near $h \sim 3\times10^{-3}$; a first
trial at the full radius $0.05$ lands deep in the fold, and the
inner solver's quadratic model, poisoned by that single rejected
step's spurious curvature, stops evaluating the objective and
returns the start point unchanged --- while $+7\times10^{4}$ N was
available at $h = 10^{-3}$.}
\label{fig:s23driver}
\end{figure}

Three defects surfaced across the campaigns, and all three are worth
recording because all three produced plausible output while being
wrong.

\textbf{No acceptance test.} The first version adopted the endpoint of
every segment unconditionally. A segment walks a local model on a
frozen record, so a long segment can walk clean past the region where
that model is valid and return a design that is genuinely worse. The
first run did exactly that: segment~1 improved the thrust, segment~2
overshot, and the carrier reported the overshoot as the answer --- a
\emph{loss} of $0.21\%$ presented as the outcome of a maximization. A
segment endpoint is a \emph{trial} point, and the fix is the standard
one: if the trial does not beat the incumbent, revert to the best
certified design and shrink the radius.

\textbf{A radius that undid its own shrinks.} With the acceptance test
in place the walk stopped losing, but it also stopped improving: it
oscillated between one good design and a succession of overshoots. The
cause was that the outer loop read its trust radius back from the
inner optimizer after each segment. That inner radius tracks the inner
method's own model, and adopting it silently overwrote every shrink
the outer loop had just ordered. Once the outer radius was kept as the
outer loop's own property --- grown only on an accepted design,
shrunk only on a rejected one --- the walk resumed and the gradient
fell by a further order of magnitude.

\textbf{The no-motion wedge.} The third defect waited for the
sharper instrument of Section~\ref{sec:gainresolved}. Started from a
converged coarse optimum at doubled resolution, the inner solver
returned its start point unchanged --- ``no motion'' --- with a
gradient of $6\times10^{6}$ in hand, and the driver, reading no
motion as convergence, stopped. The probe in
Figure~\ref{fig:s23driver}(b) exonerated the objective: smooth,
matching its linear prediction over five decades of step size,
folding over only beyond $h \sim 3\times10^{-3}$. The wedge is the
inner solver's: its first trial at the full radius lands past the
fold, the rejected step teaches its curvature model a violently
wrong lesson, and the poisoned model proposes nothing further worth
evaluating --- while $+7\times10^{4}$ N of improvement stood one
$10^{-3}$-step away. The repair stays inside the driver's own
semantics: no motion above the radius floor is treated exactly like
a rejected base --- halve the radius, retry on the same record ---
and only no motion \emph{at} the floor is convergence. Every
previously recorded walk reproduces bit-for-bit under the repaired
driver; the retries can be seen doing their work in
Figure~\ref{fig:s23driver}(a).

None of the three defects announced itself. The first produced a
confident wrong sign; the second a plausible-looking plateau easy to
report as convergence; the third a false ``converged'' verdict with
an order of magnitude of gradient still standing. One theme: the
optimizer's bookkeeping is part of the experiment, and it is graded
by the same standard as the physics.

A further gate is inherited rather than new: a design whose march
cannot be Newton-certified is rejected outright and the radius
shrunk. This fires in practice --- one segment of the S21 run of
record hit a record with a worst residual of $2.3$ and was reverted;
the purple crosses of Figure~\ref{fig:s23driver}(a) are the same
gate firing at the doubled resolution --- so it is load-bearing
rather than decorative.

\section{Does the basis contain the incumbent?}
\label{sec:reprladder}

Before any gain can be believed, one thing must be checked: can the
spline family even \emph{express} the contour it is starting from? If
not, an apparent improvement might only be the optimizer climbing back
out of a hole its own basis dug.

A spline on finitely many knots cannot reproduce the streamline
exactly, so the honest question is not whether the error is zero but
whether it \emph{falls} as knots are added. Measured on a ladder of
$6, 12, 24, 48$ knots, the largest deviation between the spline and
the streamline behaves as follows:

\begin{center}
\begin{tabular}{rccl}
\toprule
knots & max deviation & \% of spike drop & where it sits \\
\midrule
 6 & $6.536\times10^{-4}$ m & 0.0843\,\% & $x = 2.343$ m \\
12 & $4.451\times10^{-4}$ m & 0.0574\,\% & $x = 0.612$ m \\
24 & $3.418\times10^{-5}$ m & 0.0044\,\% & $x = 2.448$ m \\
48 & $2.309\times10^{-5}$ m & 0.0030\,\% & $x = 0.560$ m \\
\bottomrule
\end{tabular}
\end{center}

The total reduction across the ladder is a factor of $28$, so the
basis is dense in the sense required. But the ladder is not clean, and
its untidiness is the interesting part. From $6$ to $12$ knots the
error falls by only $1.5\times$, which for cubic interpolation of a
smooth curve is far too slow --- and taking that single ratio as the
verdict, as a first version of this check did, produces a failure
report about a basis that is in fact perfectly adequate.

Two different errors are competing, and the last column shows which
one is winning. One lives at the \emph{right} end, where the spline's
free end condition is imposed; the other sits at $x \approx 0.56$ to
$0.61$ m, the station where the streamline's curvature climbs from
essentially zero at the clamped start to its peak --- the one stretch
of the contour that genuinely bends. Neither dominates uniformly:
the location of the maximum alternates between them as knots are
added, which is exactly why no two adjacent counts can be trusted to
report a rate.

The lesson is a design lesson. Evenly spaced knots spend their budget
uniformly along a contour that is nearly straight over most of its
length, and starve the single region that curves. That is precisely
the situation the adaptive-knot machinery answers, and since this wall
uses the bell's basis, that machinery can be pointed at it without
modification. It is the natural next step and is listed as such in
Chapter~\ref{ch:status}.

One control matters here: the streamline being compared against is
itself a numerical integration, so it could in principle be the
inaccurate party. Refining its step size moves its endpoint by
$5.4\times10^{-8}$ m, nearly three orders of magnitude below the
smallest deviation on the ladder. The reference is not the
limitation.

\section{What the optimizer finds}
\label{sec:splineres}

The run of record uses six free knots on a spike of length $2.5$ m, at
the inlet angle Chapter~\ref{ch:compare} selected, in the same posed
world as every other configuration comparison.

The mass constraint is met to $4.8\times10^{-12}$ relative --- machine
precision, once the constraint is set and graded by the same
functional (Appendix~\ref{app:twoinstruments}) --- inside a derived
band of $3.8\times10^{-4}$.

The adjoint is verified next, and doing so honestly required
abandoning the obvious instrument. A finite difference cannot simply
be trusted here: the replayed march is only \emph{piecewise} smooth
in the design vector, because the wall-search topology is frozen at
the record and a perturbed design can sit fractionally on the wrong
side of a decision that record already made. Shrinking the step does
not fix this. Measured over steps of $10^{-4}$ down to $10^{-8}$, the
difference between the finite difference and the adjoint does not
fall as $h^2$; it \emph{wanders} --- for one node it runs
$-197, +61, +1.8, +0.63, -0.42$, and for another it never settles
below a fixed level at all. That is noise, not truncation, and a band
estimated from two steps describes none of it. An earlier version of
this check did exactly that and reported a correct adjoint as a
failure.

The reliable reading is the scatter itself. Taking central
differences on a ladder of steps small enough that truncation is
negligible, the spread across the ladder measures the difference's
own reliability, and that is what the adjoint may be graded against:

\begin{center}
\begin{tabular}{lccc}
\toprule
node & adjoint & difference & $|$discrepancy$|$ (band) \\
\midrule
0 & $+5.569035\times10^{6}$ & $+5.568800\times10^{6}$ &
  $2.4\times10^{2}$ $(5.2\times10^{2})$ \\
3 & $-3.523614\times10^{5}$ & $-3.523618\times10^{5}$ &
  $4.6\times10^{-1}$ $(6.3\times10^{1})$ \\
5 & $-1.119305\times10^{6}$ & $-1.119305\times10^{6}$ &
  $7.5\times10^{-4}$ $(2.1\times10^{1})$ \\
\bottomrule
\end{tabular}
\end{center}

\noindent and the directional identity $\langle \nabla J, v\rangle$
reads $9.63322\times10^{5}$ against $9.63302\times10^{5}$, a
discrepancy of $2.1\times10^{1}$ inside a band of
$1.6\times10^{2}$. Two of the three nodes agree with the adjoint far
below any plausible noise floor; the third is the one whose
difference was measured to be unreliable.

The walk then proceeds through seven accepted segments and nine
records:

\begin{center}
\begin{tabular}{lcc}
\toprule
 & $J$ [N] & $|\nabla J|$ \\
\midrule
streamline (start) & $1.15378731\times10^{8}$ & $5.844\times10^{6}$ \\
segment 1          & $1.15430463\times10^{8}$ & $1.837\times10^{6}$ \\
segment 4          & $1.15471843\times10^{8}$ & $1.681\times10^{6}$ \\
segment 5          & $1.15482085\times10^{8}$ & $2.306\times10^{5}$ \\
final              & $1.15483517\times10^{8}$ & $1.472\times10^{5}$ \\
\bottomrule
\end{tabular}
\end{center}

The gradient falls by a factor of $40$ and the thrust rises by
\begin{equation}
  \frac{J_{\text{spline}}}{J_{\text{streamline}}} - 1
  \;=\; +0.0908\,\%.
\end{equation}
In coefficient terms $C_F$ moves from $1.45163$ to $1.45295$. The
final spike descends monotonically from $1.4936$ m to $0.8330$ m and
clears the axis, so it remains a physically admissible inward-turning
contour.

\section{The gain is not resolved, and how that was found out}
\label{sec:notresolved}

The paragraph above states a gain. It does not state that the gain is
\emph{real}, and on the evidence of this carrier it is not yet
established.

An earlier version of this chapter compared the $+0.0908\%$ against a
thrust band of $0.0094\%$ and concluded it was ``about ten times the
band''. That band was inherited: it was derived in
Chapter~\ref{ch:compare} at a different resolution, on a different
computation. Quoting it here is the same error this program has now
met three times --- grading a result with an instrument that did not
produce it --- and it was flagged as a caveat before it was tested.

Testing it is cheap and the carrier now does it as check C-6: refine
the march that produced \emph{both} thrusts and let it measure its own
discretization error on the \emph{same} two designs, so whatever
cancellation the difference enjoys is preserved. On a run with ten
knots and twenty-six accepted segments:

\begin{center}
\begin{tabular}{lc}
\toprule
 & gain over the streamline \\
\midrule
at the carrier's resolution $(K, N)$      & $+0.1302\,\%$ \\
refined to $(2K-1,\; 2N-1)$               & $+0.0820\,\%$ \\
\midrule
band this refinement supports             & $0.1930\,\%$ \\
\bottomrule
\end{tabular}
\end{center}

\noindent The band is \emph{larger than the gain}. C-6 therefore
FAILS, and the honest reading is:

\begin{quote}
\textbf{It is not established that the free-form spike beats the fan
streamline. The measured gain lies inside the discretization error of
the computation that measured it.}
\end{quote}

This does not overturn the machinery, which is verified independently:
the adjoint against its finite differences, the basis against its
refinement ladder, the mass constraint to $5\times10^{-12}$, the
driver against its own rejector. Nor does it say the gain is absent
--- both evaluations are positive, and the refined one is still
$+0.082\%$. It says the experiment is not yet sharp enough to
distinguish a real gain of this size from its own noise, which is a
statement about resolution rather than about nozzles.

Two routes would sharpen it, and both are ordinary work: run at a
resolution where the band falls below the effect, and place the knots
where the contour bends (Section~\ref{sec:reprladder}) so that fewer
of them buy more shape. What should \emph{not} be done is to quote the
gain against a band borrowed from elsewhere, which is how the claim
came to look ten times stronger than it is.

\section{How much of this to believe}

The gain is small, and it should be. The exit speed of the exhaust is
fixed by energy alone, so a perfectly expanded design already sits
close to a ceiling that geometry cannot raise
(Chapter~\ref{ch:compare}). Against a near-optimal incumbent, a
recovery of one part in a thousand is the right order of magnitude; a
gain of several percent would have been evidence of a defect, not of a
better nozzle.

The measurement is also a \emph{differential} one, which is its main
protection: both thrust values come from the same march at the same
resolution, so discretization error largely cancels between them. The
band it is quoted against, however, was derived at the production
resolution of Chapter~\ref{ch:compare} and should be re-derived at
this carrier's resolution before the ratio is treated as final.

Three limitations are stated plainly. The final gradient is
$1.47\times10^{5}$, not zero, so this is a design the optimizer was
still improving when the segment budget ran out. Six knots is a small
basis, and Section~\ref{sec:reprladder} gives reason to expect that
adaptive placement would do better than more of them. And the result
stands at a single design point: one length, one inlet angle, one
ambient pressure. Above all,
Section~\ref{sec:notresolved} shows the gain is not resolved by this
carrier's own refinement, so it is a measurement awaiting a sharper
experiment rather than a result.

\section{The sharper experiment: knots bought where the contour bends}
\label{sec:adaptive}

Both routes proposed at the end of Section~\ref{sec:notresolved} have
now been run: knots placed by an error indicator rather than
uniformly, and the gain re-graded on a three-point resolution ladder
of the march that produced it. This section reports both, and the
second one is the result.

\paragraph{The indicator.} The adaptive procedure is the bell's
--- solve, mark the knot intervals that carry the residual, insert,
re-optimize, repeat --- with one substitution that the plug makes
possible. The bell reads its optimality residual off the flow, as the
drift of an adjoint-equal quantity along a control surface, because
that is the only access it has. The plug carrier owns its discrete
adjoint outright, so the same object is available directly: enrich
the basis with a candidate node at the midpoint of every knot
interval, re-interpolate the incumbent wall (for a cubic spline on a
knot superset with the same end conditions this reproduces the wall
\emph{exactly}, by uniqueness --- measured at $4\times10^{-16}$ m),
and read the gradient of thrust at the candidate nodes from one march
and one reverse-mode pass. At a class optimum the gradient at the
\emph{existing} nodes is negligible, so a candidate's component
measures precisely the freedom the current class cannot express.
The interval whose candidate carries $84.8\%$ of the total indicator
mass is the first one, $x \in [0.350, 0.565)$ --- the same
neighbourhood the representation ladder of
Section~\ref{sec:reprladder} identified from geometry alone.

\paragraph{One knot pays; a matched uniform knot does not.} Inserting
the marked candidate at $x = 0.4575$ m and re-optimizing from the
(exact) warm start raises the thrust by $+0.1921\%$ at the working
resolution. The control experiment gives that number its meaning: the
\emph{uniform} class at the same node count, optimized by the same
driver from the same kind of start, gains $+0.0002\%$ --- the extra
uniform knot lands where the contour is nearly straight and buys
nothing. Placement, not node count, is what pays. A second adaptive
insertion (at $x = 0.4037$) buys nothing further: the driver cannot
extract even one part in $10^{15}$, and the procedure stops itself
honestly. Whatever this basis can reach at this length, one
well-placed knot reaches it. Figure~\ref{fig:spikedesigns}
draws the designs as objects: the bulge the knot bought, the tail
both optimizers surrender, and the two classical machines the next
section compares against.

\paragraph{The ladder, and the verdict.} The gain of the adaptive
optimum over the fan streamline was then graded the way
Section~\ref{sec:notresolved} demands: both designs frozen, the march
refined twice, each rung re-posing its own start line, the band
derived from the ladder itself.

\begin{center}
\begin{tabular}{lcc}
\toprule
march resolution & adaptive ($m=11$) & uniform control ($m=11$) \\
\midrule
$(K, N) = (61, 51)$    & $+0.313\,\%$ & $+0.127\,\%$ \\
$(121, 101)$           & $+0.186\,\%$ & $+0.068\,\%$ \\
$(241, 201)$           & $+0.080\,\%$ & $+0.029\,\%$ \\
\midrule
band the ladder supports & $0.426\,\%$ & $0.157\,\%$ \\
\bottomrule
\end{tabular}
\end{center}

\noindent The band is still larger than the gain, at sixteen times
the cost of the original experiment. The claim that the free-form
spike beats the fan streamline remains \emph{not established} --- and
the ladder now says something sharper than ``not yet''. The measured
gain \emph{falls monotonically} under refinement, for the adaptive
design and for the uniform control alike, with a rung-to-rung decay
ratio near $0.84$. The reading: a large share of any coarse-grid gain
is the optimizer fitting the march's own discretization error ---
the design was tuned at the coarse resolution, so it is flattered by
exactly the error it was tuned against --- and the march converges
too slowly on this differential quantity for brute-force refinement
to win; each halving of the grid spacing removes only a sixth of the
remaining artifact while multiplying the cost by six.

\begin{figure}[htbp]
\centering
\includegraphics[width=\textwidth]{fig_spike_designs.pdf}
\caption{The designs this chapter compares, drawn as objects.
(a)~Three machines in one world, $x$ axial against $y$ radial [m]:
the fan-streamline spike with $\theta_E = 0$, cut at $L = 2.5$ m
(solid), with the free-form optimum of this section dashed --- the
two coincide over most of the length and separate at the tail, which
the optimizer pulls $87$ mm lower; the classical Rao construction
expanded to the same ambient (GENO, posed by exit Mach), whose spike
runs to the axis near $x = 5.9$ m --- the shaded region is the $57\%$
of it our fixed length removes; and the mass-plus-length design at
$2.5$ m, the machine of the void A/B of
Section~\ref{sec:vsrao}: with both of Rao's degrees of freedom
consumed, ambient becomes an output and its wall barely tapers.
(b)~What shape freedom actually bought at $L=2.5$ m: the adaptive
and uniform optima minus the streamline, in millimetres. \emph{Drawn
from the re-adjudication run of 2026-09-15
(Section~\ref{sec:readjudication}), not from the 08-10 record the
text of Section~\ref{sec:adaptive} describes:} the adaptive design is
now $m = 12$ with two inserted knots ($x = 0.4575$, then $0.4037$,
where the indicator put $90\%$ and $79\%$ of its mass), a $+13$ mm
bulge at $x = 0.42$ and a monotone dive to $-108$ mm at the exit; the
uniform $m = 12$ control bulges $+16$ mm at $x = 0.55$ and dives to
$-80$ mm. Both optimizers tell the same story (raise the wall where
pressure is high and large-radius, surrender radius at the tail) ---
and both, as Section~\ref{sec:readjudication} shows, do it with a
compression corner at the foot that folds the march. Thin vertical
lines mark the adaptive knot set.}
\label{fig:spikedesigns}
\end{figure}

\paragraph{What follows.} More knots are ruled out (the second
insertion bought nothing), and a longer ladder is a losing race (the
decay ratio says so). What remains is a better \emph{instrument}: an
incumbent optimized at the finer resolution, so that it is not tuned
to the coarse march's error; or an extrapolation of the paired
difference down the existing ladder, whose own band would settle
whether anything survives at the limit; or the exact two-route thrust
functional of Chapter~\ref{ch:status}, once its field residual is
understood. Until one of those runs, the honest summary of the
free-form spike is unchanged in verdict and sharpened in diagnosis:
the machinery is verified, the placement principle is demonstrated,
and the headline gain is a property of the instrument's resolution,
not yet of the nozzle.

\section{The gain, resolved}
\label{sec:gainresolved}

The better instrument has now been built and run, and the question
this chapter carried through three campaigns is closed. Three things
were done, in the order the evidence demanded.

\paragraph{Extrapolation, done per design.} Each design's thrust,
measured on a resolution ladder whose spacing halves per rung, is its
own convergent sequence, and each sequence must be extrapolated by
its own contraction --- the streamline's converges cleanly at first
order (measured ratio $0.506$), while a design tuned at a coarse
resolution is \emph{non-monotone} near its tuning rung and refuses a
one-mode fit. Extrapolated correctly, the S22 optimum's limit lands
$0.012\%$ \emph{below} the streamline's: the coarse-grid gain was not
merely unresolved, it was wholly the optimizer fitting the march's
discretization error, slightly past the point of self-harm.

\paragraph{An incumbent that is not tuned to the coarse march.} The
optimum was re-optimized at the doubled resolution $(121,101)$ ---
same class, same driver, sharper instrument. (Doing so exposed a
driver defect worth recording: from this start the trust-region
inner solver's first rejected trial, landing past the fold scale of
the replayed surface, poisoned its quadratic model and it returned
the start point unchanged while $+7\times10^{4}$ N of improvement
was measurably available at a step of $10^{-3}$. The outer driver
now applies its reject-and-shrink semantics to that case; every
previously recorded walk reproduces bit-for-bit.) The fine-tuned
optimum gains $+0.0098\%$ over the S22 design at its own resolution
--- and its ladder then repeats the pattern one rung later: whatever
resolution tunes the design, the march flatters it there and
retracts the flattery one rung above.

\paragraph{The paired ladder to $(1921,1601)$, and the verdict.}
The difference of the two designs, however, becomes a clean sequence
once the design leaves its tuning window, and it was extended to six
rungs --- affordable because a profiling pass found the march
spending its high-resolution hours inside a single degenerate
tensor-stacking compilation, which a one-line change (bit-identical,
gated on the record) reduced from hours to minutes:

Figure~\ref{fig:s23ladder} carries the result. Panel~(a) is the
whole story in one frame: the streamline's thrust climbs
monotonically onto its extrapolated limit with ratios locked at the
march's first order, while the optimum --- tuned at the second rung
--- overshoots there and then descends onto the \emph{same} limit;
the two designs are one nozzle measured two ways. Panel~(b) shows
the paired gain minus its extrapolated limit falling as a straight
line on a log scale --- geometric convergence, ratios printed at
each rung --- and the open marker is the sixth rung as the model
\emph{predicted it before it ran}: the miss is eighteen times
smaller than the step it predicted, which is the model earning its
extrapolation. That limit is

\begin{figure}[htbp]
\centering
\includegraphics[width=\textwidth]{fig_s23_ladder.pdf}
\caption{The six-rung paired ladder, \emph{drawn from the
re-adjudication run of 2026-09-16} (Section~\ref{sec:readjudication};
the 08-11 record described in the text of this section is history and
its figure was lost). (a)~Thrust of both designs against march
resolution, spacing halving per rung: the fan streamline (blue)
converges monotonically at first order (successive-difference ratios
$0.51, 0.50, 0.50, 0.50$) onto the dashed extrapolated limit; the
re-adjudicated fine optimum (orange), tuned at the second rung from a
streamline start, is flattered there by the instrument's own
discretization error and descends past the streamline's limit as the
march sharpens. (b)~The paired gain minus its extrapolated limit: a
straight line on the log scale is geometric convergence, with the
ratio settling toward the march order (printed: $1.87, 0.88, 0.64,
0.56$); the open red circle is rung six as predicted from rungs three
through five before it was computed. The extrapolated limit,
$-0.096\% \pm 0.181\%$, is zero within its band but the band exceeds
the bar ($0.040\%$): UNRESOLVED --- and, since the fine design's march
is folded, not a gain at all.}
\label{fig:s23ladder}
\end{figure}
\begin{equation}
  \text{gain}^{*} \;=\; -0.011\,\% \;\pm\; 0.017\,\%
  \qquad\text{(band from the limit's own window stability)}.
\end{equation}

\begin{quote}
\textbf{The gain is zero within the band. At this length and this
inlet split, the truncated fan streamline is optimal to within
$0.028\%$: the free-form spike does not beat it, and every gain this
program ever measured for it was the instrument's.}
\end{quote}

Rao's transposed argument (Section~\ref{sec:whynotstreamline})
survives in sign --- the measured gain is positive at every finite
resolution --- and turns out to be worth approximately nothing in
the limit, \emph{at the optimized inlet split}. The one design
variable at fixed length that pays remains the inlet angle, at
percent scale; shape refinement after it is bounded by three parts
in ten thousand. The statement is scoped: it is made at
$\theta_E = 0$ and $L = 2.5$ m, and says nothing about deliberately
suboptimal inlet splits, other lengths, or the full-expansion
comparison of Section~\ref{sec:vsrao}, which remains the open
experiment.

\section{Re-adjudication (2026-09-15/16): the optimum leaves the class}
\label{sec:readjudication}

The artifacts of the two preceding sections (the S22 adaptive design,
the S23 fine optimum, its ladder, and the figures drawn from them)
were lost with the working copy on 2026-09-04 and could only be
regenerated by re-running the carriers on the rebuilt line. That
re-run is a \emph{re-adjudication}, not a reproduction: the march in
use is the one after the general-inlet work (Chapter~\ref{ch:generalinlet}),
and the S22 driver of record had itself been changed by S23 (the
reject-and-shrink retry) before the S22 record was taken. The numbers
printed in Sections~\ref{sec:adaptive} and~\ref{sec:gainresolved} are
therefore \textbf{history}: their design vectors are gone, and the
runs of record for both carriers are now those of 2026-09-15/16. What
those runs found changes the reading of both sections, and is
recorded here in the order it was measured.

\paragraph{S22 re-run, with the ladder gated on certification.} The
adaptive carrier was amended first: every rung of its A-8/A-9 ladders
is a measurement only if \emph{both} marches certify, an uncertified
rung voids the ladder, and a void ladder is a FAIL by discipline (the
record of 08-10 had quoted rung gains without asking). Re-run at the
posing of record ($m_0 = 10$, $(61,51)$, three cycles), the carrier
reproduces itself to the bit (design and class-0 optimum identical
across two same-day runs) and produces a design one cycle richer than
the record's: two adjoint-guided knots, at $x = 0.4575$ (cycle 1,
$78.8\%$--$90.2\%$ of the indicator mass in the first interval, as
before) and then $x = 0.4037$ (cycle 2), $m = 12$, $J =
1.16199981\times10^{8}$ N at the working resolution, $+0.3975\%$ over
its class-0 start. Its ladder is VOID: the adaptive design does not
certify at $(121,101)$ (worst Newton residual $3.25\times10^{11}$, in
the edge cell of the first column) and the \emph{incumbent} certifies
only marginally at $(241,201)$ ($1.59$); the uniform control's ladder
is void at its third rung likewise ($16.7$). Verdict $11/13$, the two
FAILs being A-8 and A-9 themselves; the printed gains
($+0.521/+0.429/+0.260\%$, band $1.04\%$) are not trusted.

\paragraph{What the design is: a compression corner, and a folded
march.} The reason the ladder cannot grade these designs was found by
looking at them rather than at their thrust. The plug's wall lies
\emph{below} the flow, so at the foot a wall angle less negative than
the incoming flow angle is a compression, the opposite of the bell's
convention. The incoming flow at the start station is at
$-26.66^\circ$; the adaptive $m = 12$ design leaves it at
$-11.4^\circ$ and the cycle-1 $m = 11$ design at $-18.1^\circ$ ---
concentrated compressions of $15^\circ$ and $8.5^\circ$ at the foot
of a nozzle whose class is shock-free by construction. Downstream of
that corner the characteristics of one family cross: on the recorded
net, $59$ of $62$ columns at $(61,51)$ and $120$ of $122$ at
$(121,101)$ are non-monotone in $y$ (the streamline: $0$ and $11$,
the latter confined to the thin row-growth cells at the free jet),
and the signed area of the true net cell, normalized by its leg
lengths (the local Jacobian of the net, a sine of the angle between
the families, $\sin 2\alpha \approx 0.67$ where healthy), reaches
$-0.80$ and $-0.83$ at $(61,51)$ and $-0.98/-0.93$ at $(121,101)$,
\emph{converging} under refinement. The certification the march
carries --- the Newton residual of every cell --- is blind to this:
a folded cell solves its compatibility relations as well as any
other. The thrust these designs report is booked downstream of
$x \approx 0.75$ m through the tangled net (the first $0.4$ m
\emph{lose} thrust), a number taken from a multi-valued solution.
The $(121,101)$ certification failure sits upstream of the first
fold, in the edge cell of column 0: a lottery of the resolution,
caused by the corner rather than by the fold. (These censuses were
taken by session probes on the recorded net --- S28 addendum B and
the S29 log --- and are to be carried by the margin carrier below.)

\paragraph{S23 re-run, from a certified start.} The gain-resolution
carrier was amended to require its warm start to certify at the fine
instrument (rule R-2c), with the declared fallback to the fan
streamline on the S22 knot set. The fallback fired (both S22 designs
fail at $(121,101)$, as above); the TR-SQP walk at $(121,101)$ from
the streamline took $29$ segments to $J = 1.16287610\times10^{8}$ N,
$+0.4135\%$ over its start, certified ($0.484$), adjoint against
finite differences inside the band --- and the fine optimum is
\emph{the same object again}: wall angle at the first station
$-12.0^\circ$ at $(61,51)$ and $-14.9^\circ$ at $(121,101)$, the
design raised at the foot ($+12$ mm at the first knot) and lowered
monotonically to $-68$ mm at the tip, $59/62$ and $120/122$ folded
columns, cell margin $-0.82$ and $-0.97$ against the streamline's
$+0.52$ on the same cells. From a clean start the unconstrained
instrument reproduces the S22 mechanism: the optimizer leaves the
shock-free class as soon as it is allowed to, because the lever it
finds --- a compression corner at the foot --- is exactly the
direction the class forbids.

\paragraph{The six-rung ladder, and why its number is not a gain.}
Extended to $(1921,1601)$ as in Section~\ref{sec:gainresolved}, the
paired ladder is single-mode and settles at the march's order
(R-5$'$a: diff ratios $1.87, 0.88, 0.64, 0.56$; R-5$'$b: rung six
predicted at $-0.040\%$, measured $-0.033\%$, miss one seventh of
the step) with gains $+0.501, +0.414, +0.251, +0.108, +0.017,
-0.033\%$. Its limit is $-0.096\%$ with a band of $0.181\%$ ---
zero within band, but \textbf{UNRESOLVED} (band $> \mathrm{BAR} =
0.040\%$), and in any case a number taken from a folded march. The
re-adjudicated verdict of [X-PGRS] is therefore not a gain at all:
it is \emph{out of class}. Figure~\ref{fig:s23ladder} now draws this
ladder; Figure~\ref{fig:spikedesigns} the $m = 12$ design;
Figure~\ref{fig:workedexample}(d) the fine optimum's step.

\paragraph{The S21 optimum, re-run: the mildest design of the line is
folded too.} The free-form carrier of Section~\ref{sec:splineres} was
re-run on 2026-09-16 at its production posing ($m = 10$, $(61,51)$,
30 segments): gain $+0.1306\%$ over the streamline (record
$+0.1302\%$), $+0.0834\%$ at the doubled resolution (record
$+0.0820\%$), band $0.1889\%$ (record $0.1930\%$) --- C-6 fails as it
did, and the optimum, $J = 1.15739891\times10^{8}$ N, is the S22
class-0 start to the printed digit. Its first knot sits $+12.4$ mm
above the streamline and the wall leaves the first station at
$-24.70^\circ$ against $-26.66^\circ$: a compression corner of two
degrees, an order of magnitude milder than the adaptive designs'
--- and the march folds all the same, $60$ of $62$ columns at
$(61,51)$, $120$ of $122$ at $(121,101)$, margin $-0.67$ and $-0.73$
where the streamline reads $+0.080$ and $+0.058$. The fold scale
measured along the optimizer's own direction (the margin crosses
its floor two to three millimetres from the streamline) is the
reason: twelve millimetres on the first knot is four times it. The
three campaigns of this chapter therefore searched outside the class
from their first accepted step, and the question they asked has yet
to be posed inside it. One measurement points at where the answer
lies: the carrier's default posing ($m = 6$, first free knot at
$x = 0.708$ m, no knot near the foot; $8/8$ with its gain $+0.081\%$
above its own band of $0.015\%$) returns a design that stays
\emph{in} the class --- wall angle at the first station equal to
the flow's, no folded column at $(81,61)$ or $(161,121)$, margin
$+0.091$ against the streamline's $+0.092$ --- and its gain barely
moves under doubling ($+0.0808 \to +0.0771\%$), where the folded
designs' gains collapsed. The two levers of the free-form spike
separate by class: surrendering radius at the tail keeps the march
shock-free; raising the wall at the foot does not. (That design is
not certified at the doubled resolution, worst residual $1.84$ ---
the same resolution lottery met above; an observation, not a
resolved gain.)

\paragraph{What this changes, and what comes next.} The statement of
Section~\ref{sec:gainresolved} --- that the truncated streamline is
optimal within $0.028\%$ --- was made \emph{inside} a class the
instruments had not enforced, and the sharper verdict of 2026-09-16
is stronger in one direction and weaker in another: the free-form
spike's measured gains at every resolution were not only instrument
fit but the fitting of a shock the march cannot represent; and the
question "is the streamline optimal in the shock-free class?" is
open again, because no run of record has yet searched that class
with its boundary enforced. The program already owns the remedy and
the plug line had never used it: the constrained problem of the
master reference (Part~VI), $\max J$ subject to $g = 0$ and to a
fold margin $m(W) \ge \mu_0$ carried as a constraint with a gradient
(never a stall, never a pass-through), executable on the bell as the
margin governor [X-MGOV], with its multiplier $\mu$ the measured
price of shock-freeness and the length-constrained optimum expected
\emph{margin-active}. Its port into the plug (margin field on the
recorded net cells, KS aggregation with derived $\rho$ and floor
ladder, a nonlinear constraint in the same TR-SQP) is the next
carrier of this line; S21--S23 are declared "unconstrained class,
superseded" and will be re-run with the margin on.

\section{Against Rao's own construction}
\label{sec:vsrao}

The obvious question is whether this machinery agrees with the
classical construction it is competing with. The independent Fortran
code designs Rao spikes directly, so the comparison is available ---
and posing it turned out to be more informative than the answer.

Posed in this document's world (same gas, chamber, mass flow and lip
radius) and asked for the spike that expands to \emph{our} ambient,
the classical construction returns a design of length $5.825$ m with
area ratio $5.148$, whose exhaust leaves at
\begin{equation}
  \theta_E = 0.025^{\circ}.
\end{equation}
That number is worth pausing on. Chapter~\ref{ch:compare} swept the
inlet angle blind, knowing nothing of Rao's construction, and its
optimum landed at $\theta_E = 0$ --- the divergence argument's
prediction. The classical construction, from an entirely different
code and an entirely different method, puts the perfectly expanded
design at $0.025^{\circ}$. Our sweep resolves $\theta_E$ only to
$\pm 2^{\circ}$, so this is a consistency check rather than a
precision match, but it is a check between two independent routes and
it passes.

It also reframes this chapter's design point. The ideal spike for
these conditions is $5.825$ m; the $2.5$ m contour optimized above is
a \textbf{57\% truncation} of it. That is a heavily cut plug, which is
exactly the regime Chapter~\ref{ch:plugcycle} studies and exactly
where the streamline construction has least claim to optimality.

What could \emph{not} be done is the direct thrust comparison at a
common length. Rao's spike family has two degrees of freedom, so
fixing mass and length consumes both and the ambient becomes an
output: at $2.5$ m the classical design implies an ambient of
$1.7\times10^{6}$ Pa, against our $7.6\times10^{5}$. Its wall barely
tapers over that length ($1.630 \to 1.548 \to 1.636$ m, $\varepsilon =
2.47$) while ours falls monotonically to $0.833$ m ---
Figure~\ref{fig:spikedesigns}(a) draws both, and the full-length Rao
construction they are both being compared against. Same nominal mass,
same nominal length, different machines --- so their thrusts are not
comparable, and comparing them would have measured the mismatch in
expansion state rather than the merits of two design methods. The
comparison that would settle it is at $5.825$ m with both mass and
ambient matched, and it is listed in Chapter~\ref{ch:status}.

\paragraph{Amendment (2026-09-16).} The numbers above ($5.825$ m,
$\theta_E = 0.025^{\circ}$) belong to a posing of the classical code that
could not be reproduced on the rebuilt line; the member of record is
the axisymmetric ideal spike run in this world with $M_E$ fixed
($\theta_E = -0.020^{\circ}$, $L = 5.926$ m closing on the axis,
expanding to $0.9949$ of the nominal ambient --- a declared residual).
More importantly, the comparison at a common length has now been made,
and posing it exposed a defect of this chapter's own construction: the
lip fan that supplied both the incumbent streamline and the Cauchy data
on the cut was \emph{planar} (a simple wave, $q$ and $\theta$ functions of
the ray angle alone) while the march is axisymmetric. Past its last
ray that streamline is horizontal; at the shared full-expansion length
it sits $0.84$ m above the classical contour, and no shape freedom
downstream of the cut can repair data that are wrong on the cut. The
repair is not a forward axisymmetric fan --- with a free wall that is
not a construction, since in axisymmetric flow every $C^-$ ray reflects
on the wall and any wall is a streamline of its own field --- but the
inverse (Goursat) march: the corner relation at the lip point, exact
there and used nowhere else, and the terminal $C^-$ ray, straight and
uniform at the ambient, marched upstream ray by ray with the same
certified interior cell (its compatibility relation is symmetric in
the roles of the given and the unknown point), the wall being the
streamline traced backward from the tip. With the source switched off
this march reproduces the corner wave to second order ($7\times10^{-8}
\to 1.8\times10^{-8}$ under doubling); with it on, and the classical
contour read only afterwards, the traced wall lies at a mean $0.8$ mm
(maximum $7.8$ mm, ladder band $4.8$ mm) from the classical member over
its whole length, the mass through the terminal ray equals the target
to the tip floor's own share ($10^{-4}$), and the free-form incumbent
of the tournament is the classical contour to within the knots'
representation band with a thrust difference of $10^{-5}$ of $J$.
Finding Rao is, as it should be, a confirmation of correctness and
not a result: the tournament question --- a certified, in-class gain
over the ideal member at shared mass, ambient and length --- is posed
on that incumbent (carriers \texttt{a1\_axi\_fan.py} and
\texttt{a1\_plug\_tournament.py}, S29 log sec.~6).

What has changed structurally is nonetheless permanent. The spike is
no longer a contour the construction hands over, but one the program
computes; every stretch of the wall now carries a number saying what
it is worth; and the plug and the bell can at last be set against each
other with the same shape freedom on both sides. Chapter~\ref{ch:compare}
found the bell ahead by $0.282\%$ with an optimized bell facing a
constructed plug. Roughly a third of that margin is recovered here.
The gap narrows; on this evidence it does not close.

\section{A worked example, end to end}
\label{sec:workedexample}

Everything this chapter has described --- the posed world, the
direct march, the functional, the adjoint, the step --- can be
watched happening on one design. Figure~\ref{fig:workedexample}
follows a single complete evaluation at the chapter's working
settings ($81$ wall stations, $61$ start rows, the $11$-knot basis),
starting from the fan streamline.

\textbf{Pose (panel a).} The chamber state and the gas fix the
tables; the ambient fixes the free-boundary pressure; the lip fixes
the fan. The straight delivery line arrives at
$\theta_i = -26.7^{\circ}$ --- the turning, from the RDE's axial
discharge toward the spike, that the outer optimization of
Chapter~\ref{ch:compare} selected so that the fan completes the
turn to $\theta_E = 0$. The start radius is then \emph{not} chosen:
a root search places the wall's origin so the start line carries the
reference mass flow, measured by the same column functional every
later verdict uses. Everything upstream of the start line enters
the problem only through the Cauchy data on it; no knot can reach
it.

\textbf{March (panel b).} The direct march builds the field column
by column --- each column a wall-to-edge characteristic, gaining a
row as the jet widens --- and solves one small Newton system per
cell: $5{,}611$ cells here, every one certified to round-off. The
Mach field shows the physics the thrust integral will harvest: the
fan's expansion arriving at the wall, the free boundary holding
$p_a$, the top rows striding far downstream along near-characteristic
chords into the widening jet. One march at these settings is a few
seconds of arithmetic and about twenty of bookkeeping --- the
record of every discrete decision, which is what the adjoint will
replay.

\textbf{Score (panel c).} The thrust splits as
$J = F_{\mathrm{in}} + \int (p - p_a)\, 2\pi y\, (-\mathrm{d}y)$:
the start line delivers $F_{\mathrm{in}} = 106.4$ MN of momentum
flux, and the wall adds $9.3$ MN more --- the running integral in
the panel, $J = 115.71$ MN in all. The wall's overpressure starts at
$4.5\,p_a$ and is still $1.8\,p_a$ at the cut: the truncated spike
discards flow it never finished expanding, which is the $57\%$
truncation of Section~\ref{sec:vsrao} seen from the wall's side.
The integrand's $2\pi y$ weight is visible as the integral's
saturation --- by $x = 1.5$ m the remaining metre of slender wall
has little area left to push on.

\textbf{Gradient, and the step it justified (panel d).} One reverse
pass through the recorded schedule prices every knot
simultaneously: $+26.8$ MN per metre of raise at the first knot ---
the fat end, where overpressure and radius are both large --- decaying
and turning negative along the tail ($-1.1$ MN/m at the last knot).
The optimizer's eventual answer, overlaid, is the gradient's
instruction executed under the certification gate: raise the front
by millimetres, surrender the tail by $87$ mm. That this step ---
worth $+0.30\%$ to the coarse march that computed it --- is worth
nothing to the exact functional is the resolution of
Section~\ref{sec:gainresolved}; the example is the method working
exactly as built, on the one problem where the honest answer was
``the incumbent was already optimal.''

\begin{figure}[htbp]
\centering
\includegraphics[width=\textwidth]{fig_worked_example.pdf}
\caption{One evaluation of the plug method, end to end, from a live
run at the working settings. (a)~The posed world: lip fan (grey
rays), the mass-set start line (green), the spike wall, and the
free boundary --- whose polyline begins where the march's first
column reaches it, one long chord downstream of the start line (the
measured first-column stride of Chapter~\ref{ch:plug}).
(b)~The direct march's mesh, one point per certified cell, coloured
by Mach: columns hug the wall, and the top rows stride into the
widening jet. (c)~Wall overpressure (left axis) against the running
wall-thrust integral (right axis): the fat end does the work, and
the flow leaves the cut still $1.8\,p_a$ above ambient.
(d)~The reverse-mode adjoint at every knot (bars, symmetric log
scale) beside the step the optimizer ultimately took (red): raise
the fat end, surrender the tail --- correct at the resolution that
computed it, and worth nothing in the limit
(Section~\ref{sec:gainresolved}). \emph{Regenerated 2026-09-16: panel~(d)'s
step is the re-adjudicated fine optimum of
Section~\ref{sec:readjudication} ($+12$ mm at the first knot, $-68$
mm at the tip), a design whose march is folded.}}
\label{fig:workedexample}
\end{figure}
